The very first, entirely Hungarian-made, satellite was ``MASAT--1'', a
nanosatellite ``cubesat''. It was made mostly of aluminium with total mass
1\,kg, sides length $l = 10$\,cm and a longer communication aerial. It was
designed and prepared at the Technical University of Budapest (BME) in 2009 by
students. The launch occurred on February 13, 2012, using a Vega rocket from
the Kourou launch site --- together with several other cubesats from other
countries. It operated successfully right up to the last minutes of its
lifetime (the transmitted last data packages were captured by radio receivers
on the evening of January 9, 2015, a few hours later the nanosatellite
re-entered the atmosphere and disintegrated).

\ioaafig{2019_TH_14-f1.pdf}

The orbital altitude of MASAT--1 changed between $h_{\mathrm{min}} = 350$\,km
and $h_{\mathrm{max}} = 1450$\,km (due to its highly eccentric orbit), but in
this question assume it to be in a circular orbit 900\,km above the sea level.

The MASAT team wished to photograph their nanosatellite from the ground. Thus,
they called the staff of Baja Astronomical Observatory (South Hungary,
$\lambda_{\mathrm{B}} = 19.010843^\circ$,
$\varphi_{\mathrm{B}} = 46.180329^\circ$, $h_{\mathrm{B}} = 100$\,m) to
photograph the orbiting MASAT--1 with their telescope.

The observatory has a Ritchey-Chr\'etien type reflecting telescope with a
diameter of 50\,cm, and focal ratio of $f/8.4$. The CCD camera installed on the
telescope had a $4096 \times 4096$ chip with 9\,$\mu$m sized square pixels. The
quantum efficiency of the CCD was about 70\,\%. The practical detection limit
was about $19.5^{\mathrm{m}}$ visual band applying a
$\tau_{\mathrm{exp}} = 2$\,min exposure time. Disregard any effects of seeing.
The orbital inclination of MASAT--1 was about $i = 70^\circ$ and the direction
of the orbital motion was the same as the Earth rotation. In addition, let us
assume that the reflection area always is 100\,cm$^2$. Throughout this event,
the telescope was pointed to the local zenith, and its RA motor was tracking
the stars. Consider only light from the Sun, you can ignore light reflected
from Earth and Moon. Ignore the effects of atmospheric extinction.

\begin{description}
\item[a)] Calculate the apparent visual magnitude of this cubesat under ideal
  observational circumstances, i.e.\ when it was at the local zenith of the
  observing site (Baja, Hungary) at midnight. Omit all atmospheric effects, and
  consider the Earth to be a sphere. The albedo of a specular aluminium plate
  is $a \approx 0.70$. \hfill(10\,p)\\
  \textit{Hint:} Use a comparison of MASAT--1 with the Full Moon.
\item[b)] What was the Observatory's answer to the MASAT team; would it ever be
  possible to photograph their nanosatellite with the existing equipment at the
  observatory? \textbf{Mark ``YES'' or ``NO'' on the answer sheet.} Support
  your answer with detailed calculations. \hfill(42\,p)
\item[c)] What would have been the answer if they had taken into account the
  blurring effect of the atmosphere, the so-called seeing? In Hungary the
  typical FWHM (Full Width at Half Maximum) value of the blurred stellar images
  --- which can be approximated by a symmetric 2D Gaussian profile close to the
  zenith --- is about $3.5''$. \textbf{Mark ``YES'' or ``NO'' or ``Close to the
  limit'' on the answer sheet.} Support your answer with a short calculation.
  \hfill(8\,p)\\
  \textit{Hint:} Although the illumination of the seeing spot in the focal
  plane can be approximated by a symmetric 2D Gaussian profile, you can take it
  to be homogeneous in your calculation.
\end{description}
